% ----- Configuración del documento -----
\documentclass[12pt,a4paper]{article}

% ----- Paquetes necesarios -----
\usepackage[utf8]{inputenc}
\usepackage{graphicx} % Para insertar imágenes
\usepackage{geometry} % Para ajustar márgenes
\usepackage{hyperref} % Para insertar hiperenlaces
\usepackage{xcolor} % Para dar color a la fuente
\usepackage{amsmath} % Operaciones matemáticas
\usepackage{amssymb} % Para símbolos como checkmarks (\checkmark) y cruces (\times)
\usepackage{float} % Para 
\usepackage{parskip} % Para espaciados entre parrafos
\usepackage{subcaption} % Para el entorno subfigure

% ----- Definiciones propias -----
\definecolor{codegreen}{RGB}{0,150,0} % Color verde consola
\definecolor{MoradoIntenso}{RGB}{200, 0, 204} % Morado intenso (R:152, G:0, B:204)
\hypersetup{
    colorlinks=true,
    linkcolor=codegreen,
    filecolor=cyan,      
    urlcolor=MoradoIntenso,
} % Para que se vean azul las referencias a imágenes, en morado los links.

% ----- Declaración de márgenes -----
\geometry{
    top=2.5cm,
    bottom=2.5cm,
    left=2.5cm,
    right=2.5cm
}

% ----- Espaciados entre párrafos y sangría ----- 
\setlength{\parskip}{11pt}   % Espaciado (10 pt)
\setlength{\parindent}{14pt}  % Sangría (15 pt)

% ----- Información documento -----
\title{\textbf{Subsonic Festival - Práctica 2}}
\author{Roberto García García\\
        Javier Cumbres Tena\\
        Ismael Ramos Márquez\\
        Mario Sequedo Rodríguez\\
        Alonso Grillo Frade
        }
\date{Curso 2025-26}
\raggedbottom



%##############################################################################
%############################## INICIO DOCUMENTO ##############################
%##############################################################################

\begin{document}

\maketitle
\begin{figure}[h]
    \centering
    \includegraphics[width=0.5\textwidth]{assets/subsonicfestival.png}
\end{figure}
\newpage

% --- ÍNDICE ---
\renewcommand{\contentsname}{Índice}
\tableofcontents
\newpage



%#############################################################################
%############################## 1. INTRODUCCIÓN ##############################
%#############################################################################
\section{Introducción}

El presente documento recoge la fase de prototipado y diseño para el desarrollo de \textbf{Subsonic Festival}, una plataforma web \textit{full-stack} orientada a la gestión integral de festivales de música. Esta aplicación nace con el objetivo de centralizar y facilitar tareas críticas de la organización de eventos, tales como la promoción de artistas, la venta de entradas, la comercialización de \textit{merchandising} y la gestión del alquiler de espacios para proveedores de servicios.

El propósito principal de esta segunda entrega (Práctica 02) es diseñar las vistas para las funcionalidades descritas en la primera entrega, sin la necesidad de desarrollar las funcionalidades al uso. 

Para cumplir con estos objetivos, el contenido de este documento se organiza de la siguiente manera:

\begin{itemize}
    \item \textbf{En la Sección 2 (Esbozo de las vistas):} Se plantean los esbozos de las vistas, desarrollando los elementos que la componen y las ligaduras entre ellos. En cada vista, se presentará la funcionalidad que soporte, como se podrá acceder a ella, a que vista se navega desde ella y como.
    \item \textbf{En la Sección 3 (Árbol de navegación):} Se presenta un diagrama de estados que muestra el flujo lógico y global de navegación de la aplicación. En dicho esquema, se presentan diferentes eventos, como actuar ante ellos, además de como se tratarán los datos al realizar la navegación entre los eventos.
    \item \textbf{En la Sección 4 (Anexos):} Se adjuntan las actas de las reuniones de coordinación, evidenciando el seguimiento del proyecto.
\end{itemize}



%#############################################################################
%############################## 2. ESBOZO VISTAS #############################
%#############################################################################
\section{Esbozo de las vistas}
\subsection{Vistas}

\begin{figure}[H]
    \centering
    % Primera fila
    \begin{subfigure}{0.49\textwidth}
        \centering
        \label{fig:vista1}
        \includegraphics[width=\textwidth]{assets/Login.png}
        \caption*{Inicio Sesión}
    \end{subfigure}
    \hfill
    \begin{subfigure}{0.49\textwidth}
        \centering
        \label{fig:vista2}
        \includegraphics[width=\textwidth]{assets/Register.png}
        \caption*{Registro}
    \end{subfigure}

    \vspace{0.5cm} % Espacio vertical entre filas

    % Segunda fila
    \begin{subfigure}{0.49\textwidth}
        \centering
         \label{fig:vista3}
        \includegraphics[width=\textwidth]{assets/Home.png}
        \caption*{Home}
    \end{subfigure}
    \hfill
    \begin{subfigure}{0.49\textwidth}
        \centering
        \label{fig:vista4}
        \includegraphics[width=\textwidth]{assets/Home.png}
        \caption*{Carrito}
    \end{subfigure}

    \vspace{0.5cm}

    % Tercera fila
    \begin{subfigure}{0.49\textwidth}
        \centering
        \label{fig:vista5}
        \includegraphics[width=\textwidth]{assets/Home.png}
        \caption*{Perfil}
    \end{subfigure}
    \hfill
    \begin{subfigure}{0.49\textwidth}
        \centering
        \label{fig:vista6}
        \includegraphics[width=\textwidth]{assets/Home.png}
        \caption*{Historial de entradas}
    \end{subfigure}

    \vspace{0.5cm} % Espacio vertical entre filas

    % Cuarta fila
    \begin{subfigure}{0.49\textwidth}
        \centering
         \label{fig:vista7}
        \includegraphics[width=\textwidth]{assets/Home.png}
        \caption*{Detalles del evento}
    \end{subfigure}
    \hfill
    \begin{subfigure}{0.49\textwidth}
        \centering
        \label{fig:vista8}
        \includegraphics[width=\textwidth]{assets/Home.png}
        \caption*{Vista proveedor}
    \end{subfigure}
    \caption{Galería de vistas del sistema}
    \label{fig:vistas_sistema}
\end{figure}
\newpage

% ---------- Breve explicación de cada una de las vistas (referenciadas) ----------
\subsection{Explicación vistas}
\begin{itemize}
    \item \textbf{\textit{\hyperref[fig:vista1]{Vista 1. Inicio Sesión}}}: La página contará con una interfaz compartida para autenticación y registro. En esta pantalla se mostrarán dos secciones: en la parte derecha se situará el formulario de inicio de sesión, donde el usuario podrá introducir sus credenciales para acceder a la plataforma. Desde esta misma interfaz existirá un botón o enlace que permitirá cambiar a la vista de registro, manteniendo el mismo diseño de pantalla. 
    \item \textbf{\textit{\hyperref[fig:vista2]{Vista 2. Registro}}}: La vista de registro utilizará la misma estructura visual que la pantalla de inicio de sesión, pero en este caso el formulario mostrado será el de creación de una nueva cuenta, situado en la parte izquierda de la interfaz. El usuario podrá introducir los datos necesarios para registrarse y también tendrá disponible un botón para volver a la vista de inicio de sesión sin cambiar de página.
    \item \textbf{\textit{\hyperref[fig:vista3]{Vista 3. Home}}}: Será la página inicial de la web. En ella se mostrará un apartado con noticias recientes o destacadas, junto con una sección situada a la derecha de la interfaz donde se visualizarán los próximos eventos, permitiendo a los usuarios consultar de forma rápida la información más relevante y las actividades que se celebrarán próximamente.
    \item \textbf{\textit{\hyperref[fig:vista4]{Vista 4. Carrito}}}: En la vista del carrito se mostrarán todos los eventos añadidos por el usuario, incluyendo el nombre del evento, el precio pagado por cada entrada y una breve descripción del evento.
    
    Además, al pulsar el icono del carrito situado en la parte superior derecha, se desplegará una vista en miniatura del carrito, donde se mostrará el precio total acumulado y se ofrecerá la opción de realizar el pago directamente. 
    \item \textbf{\textit{\hyperref[fig:vista5]{Vista 5. Perfil}}}: Permite visualizar y modificar la información personal del usuario, así como gestionar su método de pago. Además, en la parte derecha de la interfaz se mostrará una sección con los próximos eventos, de forma similar a la presentada en la página principal.
    \item \textbf{\textit{\hyperref[fig:vista6]{Vista 6. Historial de entradas}}}: Vista que muestra el historial completo de las entradas adquiridas por el usuario, incluyendo el evento al que pertenecen, la cantidad de entradas compradas y el precio correspondiente.
    \item \textbf{\textit{\hyperref[fig:vista7]{Vista 7. Detalles del evento}}}: Vista detallada de un evento en la que se muestra el nombre del evento, las fechas en las que se celebra, los artistas participantes, las entradas disponibles, el horario por día y los puestos que asistirán al evento.
    \item \textbf{\textit{\hyperref[fig:vista8]{Vista 8. Vista del proveedor de servicios}}}: Vista destinada a los proveedores de servicios, que les permite consultar y adquirir espacios dentro de los distintos festivales para instalar el puesto que deseen.
\end{itemize}



%#############################################################################
%############################ 3. ÁRBOL NAVEGACIÓN ############################
%#############################################################################
\section{Árbol de navegación}

\begin{figure}[H]
    \centering
    \label{fig1:arbol_nav}
    \includegraphics[width=0.8\textwidth]{assets/arbol_navegacion.jpeg}
    \caption{Diagrama de estados del flujo de navegación}
\end{figure}
\bigskip \bigskip \bigskip 

En el \hyperref[fig1:arbol_nav]{árbol de navegación}, se muestra un diagrama de estados que representa el flujo lógico de navegación de la aplicación. A continuación, se detallan los distintos ejes de interacción.

El primero de los estados al que accedemos una vez cargado el sitio web es la página principal (\textbf{MAIN}), ya que se puede navegar por esta aún sin haber iniciado previamente sesión. 
\newpage

% ---------- Acceso y autenticación de usuario ----------
\subsection{Acceso y Autenticación de Usuario}

En caso de requerir una acción autenticada (necesaria para el uso del carrito de compras o gestión de entradas), el flujo se divide principalmente en dos ramas:
\begin{itemize}
    \item \textbf{Inicio de Sesión:} Si el usuario tiene una cuenta creada, bastará con introducir sus credenciales en el estado \textit{Iniciar sesión}. Tras la validación, será redirigido nuevamente al estado \textbf{MAIN}.
    \item \textbf{Registro de Cuenta:} En caso de no poseer cuenta, se transita al estado \textit{Añadir cuenta / Crear contraseña}. Una vez completados los campos y guardado el nuevo usuario, el sistema redirige al estado \textbf{MAIN} con la sesión ya activa.
\end{itemize}

% ---------- Descubrimiento y adquisición de eventos ----------
\subsection{Flujo de Descubrimiento y Adquisición de Eventos}
El usuario puede localizar los distintos eventos a través de dos vías:

\begin{itemize}
    \item \textbf{Exploración Pasiva:} Mediante el estado \textit{Recomendado}, que muestra contenido sugerido o recomendado (según valoraciones, fecha de proximidad, relevancia, etc).
    \item \textbf{Exploración Activa:} A través del estado \textit{Búsqueda}, donde se aplican filtros específicos que ayuden a cercar la búsqueda de eventos específicos.
\end{itemize}
Ambas vías convergen en el estado \textit{Detalles Evento}. Desde este punto, la transición \texttt{SELECT COMPRAR} redirige al usuario al módulo \textit{Métodos de Compra}, formalizando la transacción financiera.

% ---------- Flujo para cliente/usuario ----------
\subsection{Gestión de Perfil e Historial de Usuario}
Este eje está condicionado a la existencia de una sesión activa (\textit{Si está registrado}). El estado \textit{Perfil} actúa como nodo de control para:
\begin{itemize}
    \item \textbf{Configuración:} Acceso al estado \textit{Modificar Perfil} mediante la transición \texttt{SELECT AJUSTES}.
    \item \textbf{Consulta de Activos:} Transición al estado \textit{Historial}, donde el usuario puede visualizar sus entradas.
    \item \textbf{Gestión de Reembolsos:} Si una entrada cumple la condición \textit{Si entrada no consumida}, se habilita la transición \texttt{SELECT REEMBOLSAR}, que activa el estado de ejecución \textit{Action}.
\end{itemize}

% ---------- Flujo para proveedores ----------
\subsection{Búsqueda y Reserva de Espacios}
Para usuarios con rol de proveedor, se ofrece un flujo dedicado a la infraestructura:
\begin{itemize}
    \item La transición \texttt{SELECT ESPACIOS DISPONIBLES} conduce al estado \textit{Búsqueda de Espacios}.
    \item Tras seleccionar un activo, se accede a \textit{Detalles del Espacio}, culminando el flujo en el estado \textit{Reservar Espacio}.
\end{itemize}



%#############################################################################
%################################# 5. ANEXOS #################################
%#############################################################################
\section{Anexos}

\subsection{Herramientas utilizadas}
\begin{itemize}
    \item \textbf{MODELIO OPEN SOURCE: }Para el diseño del diagrama de estados propuesto para representar el flujo lógico de navegación de la \hyperref[fig1:arbol_nav]{Figura 1}, se ha optado por el uso de la herramienta \href{https://www.modelio.org/index.htm}{Modelio}, una herramienta de modelado de código abierto.
    \item \textbf{JUSTINMIND: }Para el diseño de las vistas que conforman la web, se ha optado por usar la herramienta \href{https://www.justinmind.com/es}{JustinMind}, una herramienta gratuita de diseño y prototipado para aplicaciones web y móviles.
\end{itemize}

\subsection{Actas de reunión del grupo}

\subsubsection{Acta 1}
\begin{itemize}
    \item \textbf{Fecha y hora: } 28/02/2026, 16:30 - 18:30
    \item \textbf{Asistentes: } Roberto García, Mario Sequedo, Ismael Ramos, Alonso Grillo.
    \item \textbf{Contenido tratado: } 
    \begin{itemize}
        \item [\tiny$\blacksquare$] Revisión del enunciado general y práctica 2.
        \item [\tiny$\blacksquare$] Identificación de objetivos para práctica 2.
    \end{itemize}
    \item \textbf{Decisiones tomadas: }
    \begin{itemize}
        \item [\tiny$\blacksquare$] Asignación de tareas entre los miembros del grupo.
        \item [\tiny$\blacksquare$] Decisiciones sobre implementación FrontEnd.
    \end{itemize}
\end{itemize}

\subsubsection{Acta 2}
\begin{itemize}
    \item \textbf{Fecha y hora: } 04/03/2026, 15:30 - 17:30
    \item \textbf{Asistentes: } Todos los integrantes.
    \item \textbf{Contenido tratado: }
    \begin{itemize}
        \item [\tiny$\blacksquare$] Esbozos faltantes de las vistas.
        \item [\tiny$\blacksquare$] Desarrollo del orden de navegación lógica.
        \item [\tiny$\blacksquare$] Mecanismos de desarrollo FrontEnd.
    \end{itemize}
    \item \textbf{Decisiones tomadas: }
    \begin{itemize}
        \item [\tiny$\blacksquare$] Detalles de implementación FrontEnd.
        \item [\tiny$\blacksquare$] Flujo final árbol de navegación.
    \end{itemize}
\end{itemize}

\subsubsection{Acta 3}
\begin{itemize}
    \item \textbf{Fecha y hora: } 06/03/2026, 8:30 - 10:30
    \item \textbf{Asistentes: } Ismael Ramos, Javier Cumbres, Alonso Grillo.
    \item \textbf{Contenido tratado: }
    \begin{itemize}
        \item [\tiny$\blacksquare$] Esbozo vistas.
        \item [\tiny$\blacksquare$] Documentación práctica 2.
    \end{itemize}
    \item \textbf{Decisiones tomadas: }
    \begin{itemize}
        \item [\tiny$\blacksquare$] Reparto diseño vistas.
    \end{itemize}
\end{itemize}

\subsubsection{Acta 4}
\begin{itemize}
    \item \textbf{Fecha y hora: } 09/03/2026, 12:30 - 14:00
    \item \textbf{Asistentes: } Roberto García, Javier Cumbres, Mario Sequedo.
    \item \textbf{Contenido tratado: }
    \begin{itemize}
        \item [\tiny$\blacksquare$] Paso de vistas a JustinMind
        \item [\tiny$\blacksquare$] Finalización documentación Práctica 2.
    \end{itemize}
\end{itemize}

\end{document}
%#############################################################################
%############################# FIN DEL DOCUMENTO #############################
%#############################################################################